\chapter{Desktop Interface and Workflows}
\section{Interaction Design}
The native assistant should make the task's state understandable at a glance. The main workspace combines a request area, current plan, progress timeline and result evidence. Navigation exposes task history, granted folders and configured connections. The user should be able to inspect what the system intends to do without reading raw internal protocol messages.

The interface must distinguish planning, waiting for a decision, execution and verification. A progress animation cannot imply a percentage when the amount of remaining work is unknown. Status labels such as ``Inspecting manifests'' or ``Checking the saved report'' are more useful than unsupported precision. Errors should identify the affected step and the action the user can take.

\section{Proposed Task Workspace}
Figure~\ref{fig:ui-task} is an interface concept for the planned Qt application. It shows an explicitly granted project and a report task with evidence links. The example data is illustrative. It is not a screenshot of a completed NEXOS desktop.

\reportfigure{ui-task-workspace.png}{1}{Proposed Qt task workspace showing scope, execution stages and evidence. Interface concept, not an implemented screen.}{fig:ui-task}

The request composer is accompanied by a visible scope indicator. The plan identifies meaningful actions rather than presenting a long conversation as the only record of work. A result view links to the report and shows which checks passed. If a source changes during execution, the evidence panel should state that verification failed instead of retaining a generic green completion indicator.

Keyboard navigation, visible focus and readable text are required. Controls need descriptive labels. Status must remain understandable in grayscale and should not depend solely on color. The Qt implementation should use appropriate accessibility properties and should be tested with keyboard only operation and enlarged text.

\section{Action Preview and Approval}
The approval view presents the concrete effect of a sensitive action. For a calendar request, it includes the selected account, calendar, normalized date and timezone, duration, attendee addresses and notification behavior. For a file change, it includes the destination, overwrite status and a diff or summary of the proposed change.

\reportfigure{ui-action-preview.png}{1}{Proposed calendar action preview with explicit details and a bounded approval choice. Illustrative data only.}{fig:ui-approval}

The interface binds the displayed preview to the action identifier that the broker will execute. If the engine revises an action, the old approval is invalidated and the changed details are shown. Approval and cancellation controls should have clear, distinct wording. The interface must not preapprove an action because a timer expired or because the model predicts that the user would agree.

\section{Use Case Specifications}
\subsection{Inspect a Project}
The actor is the logged in user. Preconditions are an available project root, an approved output directory and supported manifest files. The normal flow reads manifests, summarizes declared metadata, creates a new report and verifies its identity and source digests. The postcondition is a verified report with a persistent task record. Invalid syntax, unsupported input and changed source content are alternative outcomes.

\subsection{Perform a Reversible File Operation}
This planned workflow selects files within a granted root and previews destinations. The adapter rejects path escape, unexpected links and destination conflicts according to policy. A reversible operation records enough original state to support a documented undo action. Undo is a separate checked operation, not a promise that any arbitrary change can always be reversed.

\subsection{Schedule a Project Review}
This planned workflow requires a connected Google account and the relevant calendar authorization. The user supplies or confirms attendees, date, time, timezone and duration. The connector creates the approved event and checks the returned state. Provider denial, unavailable conferencing and ambiguous timeout outcomes are shown explicitly. The displayed meeting link is taken from verified provider data rather than generated by the model.

\section{First Run Experience}
First run setup should explain whether the user is in a temporary live session or an installed environment. It checks the engine, tool profiles, available storage and model readiness. The user chooses a local model configuration or an optional remote provider and grants a demonstration folder. The first task uses a prepared fixture so that the expected output is known.

The Core ISO includes the inference runtime but may need a model download before AI requests are available. The offline demonstration variant includes a selected model after licensing and resource checks. Local offline readiness must not be confused with external service readiness; a calendar connector still needs network access and account authorization.

\section{History and Failure Presentation}
The history view groups the request, plan version, attempts, approvals and evidence. A failed attempt remains visible even if a later retry succeeds. Users should be able to open a concise technical explanation without being forced to read every log line. Exported evidence should redact credentials and unnecessary private content.

Cancellation feedback must distinguish a stopped local worker from an external action whose status is still being reconciled. When work remains uncertain, the interface should say so directly. Clear incomplete states are part of trustworthy automation and are preferable to an unsupported success message.
